\documentclass{article}
\usepackage{mathtools} 
\usepackage{fontspec}
\usepackage[UTF8]{ctex}
\usepackage{amsthm}
\usepackage{mdframed}
\usepackage{xcolor}
\usepackage{amssymb}
\usepackage{amsmath}


% 定义新的带灰色背景的说明环境 zremark
\newmdtheoremenv[
  backgroundcolor=gray!10,
  % 边框与背景一致，边框线会消失
  linecolor=gray!10
]{zremark}{说明}

% 通用矩阵命令: \flexmatrix{矩阵名}{元素符号}{行数}{列数}
\newcommand{\flexmatrix}[4]{
  \[
  #1 = \begin{pmatrix}
    #2_{11}     & #2_{12}     & \cdots & #2_{1#4}   \\
    #2_{21}     & #2_{22}     & \cdots & #2_{2#4}   \\
    \vdots      & \vdots      & \ddots & \vdots     \\
    #2_{#31}    & #2_{#32}    & \cdots & #2_{#3#4}
  \end{pmatrix}
  \]
}

% 简化版命令（默认矩阵名为A，元素符号为a）: \quickmatrix{行数}{列数}
\newcommand{\quickmatrix}[2]{\flexmatrix{A}{a}{#1}{#2}}

\begin{document}
\title{习题 16.1}
\author{张志聪}
\maketitle

\section*{4}

$D$是闭域，由闭域的定义，则有开域$S$使
\begin{align*}
  D = S \cup \partial S
\end{align*}

设$P_0$为$D$的任一聚点，要证$P_0 \in D$。

反证法，假设$P_0 \notin D$，即
$P_0 \notin (S \cup \partial S)$，等价于$P_0 \notin S, P_0 \notin \partial S$.

由$P_0 \notin \partial S$，由界点定义可知，存在
$\delta_0 > 0$，使得$U(P_0; \delta_0) \cap S = \varnothing$
或$U(P_0; \delta_0) \cap S^{c} = \varnothing$.

(i)如果$U(P_0; \delta_0) \cap S = \varnothing$，
由$P_0$是聚点，有$U(P_0; \delta_0) \cap \partial S \neq \varnothing$否则会出现矛盾，
即存在$P_1 \in U(P_0; \delta_0) \cap \partial S$，
由界点定义，$U(P_1; \delta_0) \cap S \neq \varnothing$；
由于$P_0$是聚点，类似地，对任意$\epsilon < \delta_0$，
都有$P_{\epsilon} \in U(P_0; \epsilon) \cap \partial S$且
$U(P_{\epsilon}; \epsilon) \cap S \neq \varnothing$.
此时，我们有
\begin{align*}
  \rho(P_0, P_{\epsilon}) \leq \epsilon \\
  \rho(P_{\epsilon}, Q) \leq \epsilon
\end{align*}
当$\epsilon$足够小，有$Q \in U(P_{\epsilon}; \epsilon) \cap S$且
\begin{align*}
  \rho(P_0, Q) \leq \rho(P_0, P_{\epsilon}) + \rho(P_{\epsilon}, Q) \leq 2\epsilon < \delta_0
\end{align*}
于是可得
\begin{align*}
  Q \in U(P_0; \delta_0)
\end{align*}
因$Q \in S$，故$U(P_0; \delta_0) \cap S \neq \varnothing$，出现矛盾。

(ii) $U(P_0; \delta_0) \cap S^{c} = \varnothing$，于是
$U(P_0; \delta_0) \subseteq S \subseteq D$，故$P_0 \in D$，出现矛盾。

综上，任意聚点$P_0$都有$P_0 \in D$，故$D$是闭域。

\section*{13}
令$H = \{\Delta_{\alpha}\}$.

用反证法 假设定理的结论不成立，即不能用$H$中有限个开域来覆盖$D$.

$D$是一个有界闭域，因此存在一个闭正方形$S$包含它。可以把$S$等分成
四个小的闭正方形，则在这四个小闭正方形与$D$的交集中，
至少有一个小闭正方形与$D$的交集不能被$H$中有限个开域来覆盖，
记这个小闭正方形与$D$的交集为$D_1$，则$D_1 \subset D$，
且$d(D_1) < \frac{d(S)}{2}$.

再将$D_1$按上述方法等分为四个小闭正方形，同样，至少有一个
小闭正方形与$D_1$的交集不能被$H$中有限个开域来覆盖，记这个小闭正方形与$D_1$的交集为$D_2$，
则$D_2 \subset D_1 \subset D$，且$d(D_2) < \frac{d(S)}{2^2}$.

重复上述步骤并不断进行下去，则得到一个闭区域列$\{D_n\}$，它满足
\begin{align*}
  D_{n+1} \subset D_n, \, n = 1, 2, \cdots ; \\
  d_n = d(D_n), \lim\limits_{n \to \infty} \frac{d(S)}{2^n} = 0
\end{align*}
即$\{D_n\}$是闭区域套，且其中每一个闭区域都都不能被$H$中有限个开域来覆盖。

由闭区域套定理，存在唯一的点$P_0 \in D_n, n = 1, 2, \cdots$，
由于$H$是$D$的一个开覆盖，故存在开域$\Delta$，使得$P_0 \in \Delta$。
由于$\Delta$是开域，故$P_0$是内点，即存在$U(P_0; \epsilon) \subset \Delta$，
于是由16.2的推论，当$n$足够大时，有
\begin{align*}
  D_n \subset U(P_0; \epsilon) \subset \Delta
\end{align*}
这表明$D_n$只需要$H$中的一个开域来覆盖，这与构造$D_n$时的假设
“不能被$H$中有限个开域来覆盖”相矛盾。

\end{document}